\chapter{Context, Presentation, and Needs Analysis}\label{chap:context}

\section{Introduction}
This chapter sets the context of the Riaya project. It introduces the host company, Pengguin, then presents the Riaya product and its actors: the patient, the doctor with an optional secretary, and the administrator. It describes the business model, then lists the functional and non-functional requirements, and compares Riaya with existing booking solutions.

\section{Presentation of the Host Company: Pengguin}

Pengguin is a company founded by John Jang, and its core business is language e-learning. Its main product is SpeakDuo, a platform where people learn a language by speaking face to face with a partner.

Besides SpeakDuo, Pengguin sometimes takes on web development work for outside clients. This internship was part of one of those projects: Riaya, a doctor appointment platform for Tunisia, built for a client.

\begin{figure}[H]
    \centering
    \includegraphics[width=0.45\textwidth]{images/pengguin_logo.png}
    \caption{Pengguin Company Logo}
    \label{fig:pengguin_logo}
\end{figure}

\section{Presentation of the Riaya Product (Client of Pengguin)}

\subsection{Definition and Genesis of the Project}
Riaya is an AI doctor appointment platform made for Tunisia. The project started from the observation that many patients are not comfortable with web portals or mobile applications, while almost everyone can make a phone call.

Patients book an appointment by calling a phone line and talking, without installing an application or creating an account. Doctors use a web tool to run their practice: their agenda, their patients and medical files, and their invoices. The phone line can only book doctors who keep their agenda up to date in that tool, so each side depends on the other.

\subsection{Value Proposition: The ``AI-first'' and ``Voice-first'' Concept}
``Voice-first'' means the phone call is the main channel for patients. We chose this on purpose: in Tunisia, a phone call is a familiar and trusted way to reach a cabinet. WhatsApp, already widely used in Tunisia, is the second channel, used after the call for confirmations, medical documents, invoices and rating links. Patients never need a website.

``AI-first'' means the patient side is built around an AI agent instead of menus or forms. The voice agent uses the Gemini Live API with the Gemini 3.1 Flash Live preview model. It listens and answers with its own voice, so the conversation flows like a normal call and the patient can interrupt it. We picked it after comparing several realtime models and chains of separate speech-to-text, language and text-to-speech models; \chapref{chap:technical} explains that choice. During the call the agent uses tools in the Riaya system to look up places, search for doctors and free slots, book or cancel requests and save the caller's name and language. It keeps its replies short, asks one question at a time and never gives medical advice.

The doctor side uses AI as well. A doctor can record a consultation in the browser, get a transcript that separates the speakers, and ask an AI assistant about medical questions, the product or the practice's own data. The assistant can prepare an invoice or a booking decision, but nothing changes until the doctor confirms it.

\subsection{The Actors of the Platform}
The platform has five actors: the patient, the visitor, the doctor, the secretary who works on behalf of a doctor, and the administrator.
\begin{itemize}
    \item \textbf{Patient:} has no account and uses only the phone and WhatsApp. The calling number identifies the line. The patient can book standard or urgent appointments for themselves or for a family member on the same phone, list or cancel pending bookings, and rate the doctor after a visit. The AI assistant handles the rest of the call, such as the language, the callbacks and the WhatsApp messages.
    \item \textbf{Visitor:} someone who has not signed up yet, usually a doctor discovering the platform. The visitor can only sign up, with an email code, and then becomes a doctor who applies for verification.
    \item \textbf{Doctor:} signs up on the website, submits professional documents and waits for an administrator's approval. Once verified, the doctor sets weekly availability, accepts or refuses phone bookings, manages patients, medical files and invoices, records consultations, uses the AI assistant, links their own WhatsApp number and manages the subscription.
    \item \textbf{Secretary:} is invited by email by a doctor on the Pro plan. A secretary handles appointments, patients, invoices and availability for that doctor, without access to medical files, recordings, the AI assistant or billing.
    \item \textbf{Administrator:} runs the platform from a back-office. The administrator verifies doctors, monitors live calls, reviews calls reported as abusive, blocks numbers, runs best-fit searches, reads the patients' ratings of each doctor and can stop sending new patients to a doctor, manages specialities and users, links the platform WhatsApp account and reads platform statistics.
\end{itemize}

\section{Business Model}

\subsection{Free and Pro Plans}
Riaya offers doctors two plans, Free and Pro. Both give the main practice features: patients, appointments, availability, invoices and medical files. Phone booking by the AI assistant is unlimited on both plans, so a patient is never turned away because of the plan a doctor chose. The plans differ on the features that cost Riaya money to run: WhatsApp messages started by the doctor, recording and transcription time, and the AI assistant. \xref{tab:plans} compares the two plans.

Doctors can upgrade at any time. Payments go through the Konnect.network API. Konnect has no built-in subscriptions, so Riaya manages the cycle itself: a daily BullMQ job at 08:00 UTC finds Pro subscriptions that end within three days and creates the next invoice with a payment link. The doctor pays it from the Subscription page to keep Pro. A Pro plan that is not renewed falls back to Free.

\begin{table}[H]
    \centering
    \caption{Comparison of Free and Pro Subscription Plans}
    \label{tab:plans}
    \small
    \begin{tabular}{>{\raggedright\arraybackslash}p{\dimexpr 0.4074\dimexpr\textwidth-6\tabcolsep\relax\relax}>{\raggedright\arraybackslash}p{\dimexpr 0.2593\dimexpr\textwidth-6\tabcolsep\relax\relax}>{\raggedright\arraybackslash}p{\dimexpr 0.3333\dimexpr\textwidth-6\tabcolsep\relax\relax}}
        \toprule
        \textbf{Feature} & \textbf{Free} & \textbf{Pro} \\
        \midrule
        Price & 0 TND & 49 TND per month \\
        Patients, appointments, availability, invoices, medical files & Yes & Yes \\
        AI phone bookings & Unlimited & Unlimited \\
        WhatsApp sends started by the doctor & 30 per month & Unlimited \\
        Automatic WhatsApp messages to patients & Not counted & Not counted \\
        Recording and transcription time & 2 hours per month & 100 hours per month \\
        Extra recording hour packs & No & 25 h for 15 TND, 100 h for 50 TND \\
        AI assistant allowance & About 20 questions per month & About 2,000 questions per month \\
        AI consultation note on each recording & No & Yes \\
        Secretary accounts & None & Unlimited \\
        Logo and stamp on invoices & No & Yes \\
        Support & Standard & Priority \\
        Payment & None & Monthly invoice via Konnect.network \\
        \bottomrule
    \end{tabular}
    \\[0.2cm]
\end{table}

\subsection{Plan Limits}\label{subsec:usage-limits}

The two plans set their limits for different reasons. The Free limits give a doctor enough credit to try each feature, but they are too small to run a practice on, so a doctor who wants to use these features in daily work subscribes to Pro. The Pro limits follow what each feature costs us to run.

Automatic WhatsApp messages to patients, such as confirmations and review requests, have no limit on either plan. Recording time is counted when a recording is transcribed. The AI assistant is limited by what its replies cost us and not by the number of messages, because a long chat with a transcript attached costs far more than a short question. All counters reset every month.

\section{Functional Requirements}

\subsection{Functional Requirements of the Patient}
The patient has no account and uses only the phone and WhatsApp. The patient must be able to:
\begin{itemize}
    \item \textbf{Talk to the assistant in their language:} the call can be held in Tunisian Derja, French or English.
    \item \textbf{Book an appointment:} the patient asks for a nearby doctor of the right speciality, or names a doctor directly, and chooses a free slot.
    \item \textbf{Book for a family member:} one phone can serve a whole family, so the patient can book for a relative who uses the same phone.
    \item \textbf{Send an urgent request:} in an urgent case, one request goes to several nearby doctors at once and the first one who accepts keeps it.
    \item \textbf{Manage pending bookings:} in a later call, the patient can list the bookings still waiting for a doctor and cancel them.
    \item \textbf{Rate the doctor:} after the visit, the patient receives a link to rate the doctor and leave a comment.
\end{itemize}

\subsection{Functional Requirements of the Doctor}
The doctor must be able to:
\begin{itemize}
    \item \textbf{Sign up and get verified:} the doctor creates an account, submits professional documents and waits for an administrator's approval.
    \item \textbf{Manage availability and appointments:} the doctor sets the weekly working hours and manages appointments on a calendar.
    \item \textbf{Answer phone bookings:} the doctor accepts or refuses each booking made over the phone, and sees what the caller said.
    \item \textbf{Manage patients and medical files:} the doctor keeps a file for each patient, with prescriptions, consultation notes and lab reports.
    \item \textbf{Create invoices:} the doctor bills a patient, records payments and sends the invoice on WhatsApp.
    \item \textbf{Record consultations:} the doctor records a consultation in the browser and gets a transcript that keeps the doctor's words apart from the patient's. Riaya then writes a structured consultation note that the doctor can edit or regenerate, and a recording attached to a patient appears in that patient's file. Recording time counts against the plan, and extra hours can be bought.
    \item \textbf{Use the AI assistant:} the doctor asks questions in a chat about medicine, about Riaya or about a recorded consultation. The assistant reads the practice's data (schedule, pending bookings, patients, free slots, invoices and statistics) and can prepare an invoice, a booking decision, a medical note, a new appointment or a new time for an existing one. Each change appears as a card, and nothing is saved until the doctor accepts it.
    \item \textbf{Link a WhatsApp number:} the doctor links a personal WhatsApp number, so that documents and invoices reach patients from that number.
    \item \textbf{Manage the subscription and the team:} the doctor chooses a plan and, on the Pro plan, invites secretaries.
\end{itemize}

\subsection{Functional Requirements of the Secretary}
A secretary is invited by a doctor on the Pro plan. The secretary must be able to:
\begin{itemize}
    \item \textbf{Manage the doctor's appointments:} create and move appointments on the doctor's calendar, and accept or refuse phone bookings.
    \item \textbf{Manage patients and invoices:} add patients, create invoices and record payments for that doctor.
    \item \textbf{Manage availability:} set the doctor's weekly working hours.
\end{itemize}
The secretary has no access to medical files, recordings, the AI assistant or billing.

\subsection{Functional Requirements of the Administrator}
The administrator must be able to:
\begin{itemize}
    \item \textbf{Verify doctors:} review each application with its documents, then approve, reject or ban the doctor.
    \item \textbf{Monitor calls:} follow live calls with their transcript and listen to call recordings.
    \item \textbf{Handle abuse:} review calls reported as abusive and block the number for a period or for good.
    \item \textbf{Search for doctors:} run the same doctor search as the phone assistant.
    \item \textbf{Follow doctor quality:} read the patients' ratings and comments for each doctor, and stop or resume referrals to a doctor with a reason the doctor can read.
    \item \textbf{Manage the platform:} manage users, specialities and the platform WhatsApp account, and read platform statistics.
\end{itemize}

\section{Non-Functional Requirements}

\subsection{Performance and Real-Time}
The assistant must answer fast enough for a natural phone conversation, and the doctor search must finish while the caller waits on the line. New bookings must appear on the doctor's calendar without reloading the page.

\subsection{Security and Confidentiality}
Each user must only reach the data of their own role and practice, and secretaries must never see medical data. Patient data must not be kept or used for training by the AI providers. The platform must respect the Tunisian data protection law (INPDP) and, where applicable, the principles of the GDPR.

\subsection{Scalability and High Availability}
The platform must handle several calls at once, and the web application and the voice server must be deployable separately, so that updating one does not stop the other.

\subsection{Resilience and Fault Tolerance}
A failure of an outside service must never leave a patient without an answer, and no plan limit may block a patient's booking.

\subsection{Ergonomics and User Experience (UX/UI)}
Patients must not need to install anything or create an account. The assistant must speak briefly and ask one question at a time. The doctor dashboard must work on mobile and desktop.

\subsection{Technical and Regulatory Constraints}
The assistant must speak Tunisian Derja and never another Arabic dialect. Place search is limited to Tunisia, patients speak in Tunisia time, and amounts are in Tunisian dinars.

\section{Existing Solutions}
Before designing Riaya, we looked at the existing ways to book a medical appointment. They fall into three groups: international web and mobile platforms, local Tunisian directories, and classic Interactive Voice Response (IVR) systems. For each one we noted its features, its strengths and what it leaves out for Tunisian patients.

\subsection{Solution 1: Doctolib (Classic Web/Mobile Platform)}
Doctolib is one of the main European platforms for online medical appointments. It was founded in France in 2013 and is widely used in France, Germany and Italy.

\paragraph{Features.} For patients, Doctolib offers doctor search by speciality, location and availability, online booking, reminders by email and SMS, access to medical documents and teleconsultation. For doctors, it provides a practice management dashboard with a calendar, patient records, billing and connections to other software. In some countries it also connects to national health insurance and electronic health record systems.

\paragraph{Strengths.} Doctolib is easy to use, has a large network of healthcare professionals and is integrated with the health systems of the countries where it operates. In France it has become the usual way to book a doctor online.

\paragraph{Limitations.} Doctolib is designed for European healthcare systems. Patients must create an account and use the website or the mobile application, which leaves out people who are not comfortable with apps or who have unreliable internet access. There is no voice assistant for patients who cannot or do not want to use the app, and no Tunisian Arabic. It does not use WhatsApp for follow-up messages, and it has no flow that sends an urgent request to several doctors at once.

\begin{table}[H]
    \centering
    \caption{Doctolib: Main Features and Limitations}
    \label{tab:doctolib}
    \small
    \begin{tabular}{>{\raggedright\arraybackslash}p{\dimexpr 0.3571\dimexpr\textwidth-4\tabcolsep\relax\relax}>{\raggedright\arraybackslash}p{\dimexpr 0.6429\dimexpr\textwidth-4\tabcolsep\relax\relax}}
        \toprule
        \textbf{Aspect} & \textbf{Description} \\
        \midrule
        Main features & Doctor search, online booking, reminders, teleconsultation, practice management dashboard. \\
        Target market & France, Germany, Italy and other European countries. \\
        Strengths & Large network, easy interface, integration with health systems. \\
        Limitations & Patient account required, no voice assistant, no Tunisian Arabic, no WhatsApp, no multi-doctor urgent booking, built for European systems. \\
        \bottomrule
    \end{tabular}
    \\[0.2cm]
\end{table}

\subsection{Solution 2: TunisieDocteur / Mondocteur (Local Directories)}
TunisieDocteur and Mondocteur are two well-known Tunisian websites for finding a doctor. They work as directories that list doctors by speciality and city, and patients can request an appointment through them.

\paragraph{Features.} These platforms offer a searchable list of doctors with their speciality, address and contact details, simple filters by speciality, city and name, and appointment requests that the doctor confirms or declines. Some of them also show patient reviews.

\paragraph{Strengths.} They know the Tunisian market and list local doctors. Patients who are comfortable with websites find them fairly easy to use.

\paragraph{Limitations.} The patient needs internet access and, to book, an account. All interaction is text on a web page, with no voice assistant. Search relies on simple filters instead of ranking doctors by distance and the nearest free slot, and there is no way to send an urgent request to several doctors at once. Follow-up does not go through WhatsApp, and while some of these sites are available in French and Arabic, none offers a spoken conversation in Derja.

\begin{table}[H]
    \centering
    \caption{TunisieDocteur / Mondocteur: Main Features and Limitations}
    \label{tab:tunisiedoctor}
    \small
    \begin{tabular}{>{\raggedright\arraybackslash}p{\dimexpr 0.3571\dimexpr\textwidth-4\tabcolsep\relax\relax}>{\raggedright\arraybackslash}p{\dimexpr 0.6429\dimexpr\textwidth-4\tabcolsep\relax\relax}}
        \toprule
        \textbf{Aspect} & \textbf{Description} \\
        \midrule
        Main features & Doctor directory, appointment requests, patient reviews, basic search filters. \\
        Target market & Tunisia. \\
        Strengths & Local focus, knowledge of the Tunisian market, listing of local doctors. \\
        Limitations & Internet and account required, no voice assistant, no urgent booking, no WhatsApp, filter search without distance and slot ranking. \\
        \bottomrule
    \end{tabular}
    \\[0.2cm]
\end{table}

\subsection{Solution 3: Classic IVR Systems (Interactive Voice Response)}
Some hospitals and clinics handle appointment calls with a classic IVR system, an automated phone menu where the caller picks options with the keypad or short spoken commands.

\paragraph{Features.} A classic IVR plays a menu (``press 1 for appointments, press 2 for information''), reads the keys the caller presses, and routes the call to a department or a human operator. Some systems also let the caller choose a date and time with the keypad.

\paragraph{Strengths.} IVR systems are simple and cheap, and they handle many calls without staff. Anyone with a phone can use them, whatever their internet access or digital skills.

\paragraph{Limitations.} The caller has to go through a fixed menu, which is slow and often frustrating, and long menus lead people to hang up. The system does not understand speech, so it cannot take a description of symptoms. It does not know where the caller is, cannot find the nearest available doctor and has no urgent flow. It sends no WhatsApp follow-up, and most systems offer one or two languages at best.

\begin{table}[H]
    \centering
    \caption{Classic IVR Systems: Main Features and Limitations}
    \label{tab:ivr}
    \small
    \begin{tabular}{>{\raggedright\arraybackslash}p{\dimexpr 0.3571\dimexpr\textwidth-4\tabcolsep\relax\relax}>{\raggedright\arraybackslash}p{\dimexpr 0.6429\dimexpr\textwidth-4\tabcolsep\relax\relax}}
        \toprule
        \textbf{Aspect} & \textbf{Description} \\
        \midrule
        Main features & Automated menu, keypad input, basic routing, appointment scheduling. \\
        Target market & Hospitals and clinics in Tunisia and elsewhere. \\
        Strengths & Simple, low cost, usable by anyone with a phone. \\
        Limitations & Rigid menu, no speech understanding, no location, no matching, no urgent booking, no WhatsApp, few languages. \\
        \bottomrule
    \end{tabular}
    \\[0.2cm]
\end{table}

\subsection{What Riaya Adds}
None of these solutions fits the way most Tunisians look for a doctor. Doctolib is built for European health systems and requires an account. The local directories know the market, but they still expect a website, an account and an internet connection. Classic IVR systems reach anyone with a phone, but they cannot understand what the caller says or find them a doctor. Riaya keeps the phone call, which every patient already knows how to make, and puts an assistant on the line that speaks Derja, French and English. The patient explains what they need in their own words, and the assistant finds nearby doctors ranked by distance and the nearest free slot, with no account and no internet needed. When the need is urgent, it sends the request to up to three doctors at once, and the follow-up arrives on WhatsApp, which most patients already use every day. For patients who are not at ease with apps or who have no reliable internet, booking a doctor then takes a single phone call.

\section{Conclusion}
This chapter presented the host company Pengguin and the Riaya product, its actors and its plans, with the way their limits are counted. It listed the functional requirements of patients, doctors, secretaries and administrators, and the non-functional requirements on speed, security, scalability, resilience, usability and constraints. The comparison with Doctolib, local directories and classic IVR systems showed what Riaya adds for Tunisian patients. The next chapter covers the project management method, Kanban, and the modeling diagrams.
